\section{Worked problems: padded batches, what graphs buy, and when fusion pays}

\subsection{The price of padding}
\textbf{Problem.} vLLM captures CUDA graphs on an H100 at batch sizes of 1, 2
and 4 tokens, every multiple of 8 below 256 and every multiple of 16 up to 512,
and pads each step to the next captured size. Steps carry 17, 100, 300 and 600
tokens --- decode tokens, or decode tokens mixed with prefill chunks. At what
size does each run, and what share of its slots is padding? Answer again for
graphs at the powers of two. When do the padded slots cost time?

\textbf{Solution.} With vLLM's sizes, 17 tokens run at 24 (29\% padding), 100
at 104 (3.8\%) and 300 at 304 (1.3\%); 600 exceed the largest graph and run
without one, paying their launches but no padding. With powers of two, 17 run
at 32 (47\%), 100 at 128 (22\%) and 300 at 512 (41\%) (all derived). A padded
slot adds arithmetic --- one more row through every matrix product --- but no
weight reads. A decode step of 17 or 100 sequences is far from compute-bound
(decode at 4,096 tokens of context stays bandwidth-bound at every batch that
fits in an H100, \Chref{napkin-math}), so its padded rows cost almost nothing.
A step of 300 tokens that is mostly a prefill chunk is compute-bound
(\Chref{prefill-decode}), and there every padded slot costs its share: vLLM's
1.3\% costs about 1.3\%, while running 300 tokens at 512 would make the step
about 1.7 times as long for no more output. Dense sizes matter where steps are
compute-bound, and the large token counts of mixed batches are exactly where
vLLM's rule is dense.

\subsection{What graphs buy, and what they leave}
\textbf{Problem.} A Llama-1B forward pass on an H100 at batch one launches
about 100 kernels and streams 2.5~GB of BF16 weights. Launches cost
2.1~$\mu$s each, or 1.3~$\mu$s replayed from a CUDA graph
\cite{spector2025megakernel}. (a) Predict the step time with and without
graphs if the kernels themselves streamed at the full 3.35~TB/s. (b) Hazy
Research measured engines of the time at no more than half the H100's
bandwidth and their megakernel at 78\%. How large is the gap between an engine
with graphs and the megakernel, and how much of it do launches explain?

\textbf{Solution.} (a) Streaming takes $2.5 / 3.35 = 0.75$~ms. Eager launches add
0.21~ms, for 0.96~ms; graphs add 0.13~ms, for 0.88~ms --- 8\% faster (derived).
(b) At half the bandwidth, streaming takes 1.49~ms, and with graphs the step is
about 1.62~ms; at 78\% it takes 0.96~ms, the megakernel's ``under a
millisecond'' (derived). The gap is about 0.66~ms, and the graph's launches,
0.13~ms, are a fifth of it. The other four fifths are bandwidth that the
separate kernels leave unused --- by Hazy Research's account, at their
boundaries: multiprocessors idle while a kernel's last wave drains, and every
kernel stalling while its first weights arrive. Graphs remove the issue cost
of a launch; the larger prize is the gap cost, and only removing the
boundaries, by fusion or a persistent kernel, collects it.

\subsection{When does fusion pay?}
\textbf{Problem.} In the decode step that opened the chapter, llama.cpp fused
each of Llama~3.2~3B's 57 normalizations with the multiplication by its weight,
on activations of 3,072 32-bit floats per token. (a) How many bytes and
dispatches does that save per decoded token, against the 2.0~GB of weights the
step reads? (b) How many bytes does it save in a 512-token prefill, which the
M1 runs at 248 tokens per second (record \texttt{2026-09-23-m1-llama32-3b})? (c)
What kind of fusion would be worth more?

\textbf{Solution.} (a) Unfused, the normalization writes its output and the
multiplication reads it back: $2 \times 3{,}072 \times 4 = 24.6$~KB per
fusion, 1.4~MB per token --- 0.07\% of the weights --- and 57 dispatches of
about 1.3~$\mu$s each, 0.07~ms of a 50~ms step (derived). At decode the
fusion's value is in the dispatches, and even they are small on this machine.
(b) A 512-token batch saves $512 \times 24.6$~KB $= 12.6$~MB per fusion,
0.72~GB in all; at the M1's measured 46--50~GB/s that is about 15~ms of a
2.07~s prefill, 0.7\% (derived). (c) Fusion pays in proportion to the traffic
of the intermediate it removes, measured against the step's total. A
normalization's output is one vector per token, while a projection reads a
whole matrix; the fusion that changed inference is the one whose intermediate
grows fastest --- FlashAttention's, which never writes the score matrix, whose
size grows with the square of the context (\Chref{flashattention}). Fusions of
a matrix product's epilogue, such as dequantizing weights inside the product
(\Chref{weight-quantization}) or quantizing an activation as the previous
kernel writes it, follow the same rule: they remove a full pass over a large
tensor.
